% ============================================================
% Capítulo 02: Técnicas básicas de recuperación de productos
% ============================================================

% Azúcares (totales y reductores) y producto

% Consejo: Este capítulo se enfoca en el Downstream. Explica cómo separar y purificar el producto de interés.

\section{Etapas del Downstream}

El \textbf{Downstream} o etapa de recuperación de productos es la fase final de un bioproceso y tiene como objetivo aislar, concentrar y purificar el producto de interés (metabolito) a partir del caldo de cultivo.

\begin{enumerate}
    \item \textbf{Separación mecánica de células:} Se eliminan las células y otros sólidos del líquido. Técnicas comunes incluyen:
        \begin{itemize}
            \item \textbf{Filtración:} Separación por tamaño usando membranas o filtros.
            \item \textbf{Centrifugación:} Separación por densidad usando fuerza centrífuga.
        \end{itemize}
    \item \textbf{Lisis celular:} Si el producto es intracelular, es necesario romper la pared y membrana celular para liberarlo. Técnicas:
        \begin{itemize}
            \item \textbf{Mecánicas:} Homogeneización a alta presión, molienda con perlas.
            \item \textbf{Químicas:} Uso de detergentes o solventes.
            \item \textbf{Enzimáticas:} Uso de enzimas como lisozima.
            \item \textbf{Físicas:} Choque osmótico, congelación-descongelación, ultrasonido.
        \end{itemize}
    \item \textbf{Concentración del producto:} Se reduce el volumen de la solución que contiene el producto. Técnicas:
        \begin{itemize}
            \item \textbf{Ultrafiltración:} Filtración con membranas de tamaño de poro específico para retener el producto y dejar pasar solvente.
            \item \textbf{Evaporación:} Eliminación del solvente por calor.
            \item \textbf{Extracción líquido-líquido:} Transferencia del producto a un solvente inmiscible.
        \end{itemize}
    \item \textbf{Purificación del producto:} Se busca obtener el producto con el mayor grado de pureza posible. Técnicas:
        \begin{itemize}
            \item \textbf{Cromatografía:} Separación basada en diferencias de afinidad, tamaño, carga o hidrofobicidad.
            \item \textbf{Precipitación:} Separación del producto de la solución añadiendo sales, solventes o cambiando el pH.
            \item \textbf{Cristalización:} Formación de cristales puros del producto.
        \end{itemize}
    \item \textbf{Acondicionamiento del producto:} Etapa final para darle la forma de presentación deseada (polvo, líquido, etc.).
\end{enumerate}

% Consejo: La estrategia de recuperación depende del tipo de metabolito (primario/secundario) y de su localización (intracelular/extracelular).

\section{Recuperación según el tipo de metabolito}

\begin{itemize}
    \item \textbf{Metabolitos extracelulares:} Son secretados al medio de cultivo. La recuperación es más sencilla, ya que no se requiere la lisis celular. Las etapas principales son la separación de células (filtración/centrifugación) y la purificación del producto desde el medio líquido.
    \item \textbf{Metabolitos intracelulares:} Permanecen dentro de la célula. Requieren una etapa de lisis celular para liberarlos, seguida de las etapas de concentración y purificación. Esto hace el proceso más complejo y costoso.
\end{itemize}

% Consejo: Para bioprocesos con células vegetales o animales, la lisis celular puede ser más delicada.

\subsection{Bioprocesos con células vegetales, animales y enzimas}

\begin{itemize}
    \item \textbf{Células vegetales y animales:} Son más frágiles que las bacterianas. La lisis celular se puede lograr con métodos suaves como la homogenización.
    \item \textbf{Enzimas:} Al ser proteínas, son sensibles a condiciones extremas de temperatura y pH. Se purifican típicamente por cromatografía de afinidad.
\end{itemize}